\chapter{Introduzione}
\subsection*{Contesto generale }
Nell’attuale panorama globale, l’efficienza dei trasporti è gestita maggiormente da sistemi informatici. Più precisamente, negli ultimi anni si è effettuata un’analisi sul trasporto pubblico, che ci ha permesso di osservare quanto quest’ultimo ha subito una trasformazione radicale. Questa trasformazione è stata accelerata dallo sviluppo di piattaforme di infomobilità, come Google Maps, dalla contemporanea evoluzione dei dispositivi mobili, che hanno reso l’accesso alle informazioni di viaggio immediato.
\\ \\
In questo contesto si inserisce il paradigma del \textit{Mobility as a Service} (MaaS). Tale paradigma rappresenta un concetto globale di mobilità che prevede l’integrazione di servizi di trasporto pubblico e privato accessibili grazie ad un unico canale digitale. \\
Attraverso piattaforme dedicate vi è l'aggregazione di diverse funzionalità: dal trasporto pubblico di linea ai servizi di sharing (car, bike) e taxi, l’utente è in grado di pianificare, prenotare e corrispondere il pagamento delle varie soluzioni di viaggio, semplificando i vari spostamenti di urbani attraverso un’interfaccia unificata.
\\ \\
Il punto di partenza di tale lavoro è l'analisi del comune di Scicli, le cui caratteristiche morfologiche e infrastrutturali richiedono soluzioni di trasporto sempre più efficienti. 
\\ \\
Nello specifico, la conformazione territoriale del comune di Scicli è articolata nel seguente modo: un centro storico barocco (sito UNESCO) e diverse frazioni marittime (Donnalucata, Cava d’Aliga, Sampieri, ecc.), situate a una distanza variabile tra i 7 e i 10 km. L’esigenza primaria è dunque la creazione di una connessione combinata tra le linee di trasporto urbano ed extraurbano, i servizi taxi e le navette turistiche stagionali. \\
Questa frammentazione geografica, unita alla forte stagionalità dei flussi turistici, rende l'adozione di una piattaforma centralizzata lo strumento ideale per colmare i gap di collegamento logistico.

\subsection*{Obiettivi della tesi }
Il seguente lavoro di tesi si pone l’obiettivo di analizzare e progettare un’architettura software basata sui principi del MaaS, configurata per la rete di trasporti della città di Scicli. Il progetto si articola in varie fasi. 
\\
Una prima fase di analisi sviluppata secondo i principi dell’ingegneria del software, e in una successiva fase implementativa, che prevede lo sviluppo di un’applicazione mobile mediante il framework \textit{React Native}.
\\
L’applicazione si propone di digitalizzare e ottimizzare l’offerta di mobilità attraverso le seguenti funzionalità chiave:
\begin{itemize}
    \item Digitalizzazione delle linee urbane.
    \item Servizi \textit{Demand Responsive Transport} (DRT), ovvero trasporto su richiesta, corse ``a chiamata'' per collegare le frazioni al centro urbano durante specifiche fasce orarie. Tale approccio mira a superare l’inefficienza dei bus sottoutilizzati nel periodo invernale, riducendo sensibilmente l’impatto ambientale e le emissioni di gas.
    \item Digitalizzazione del sistema dei pagamenti, attraverso un’interfaccia centralizzata di controllo e transazione, dei pagamenti per le aree di sosta tariffata (strisce blu).
    \item Monitoraggio in tempo reale circa la chiusura delle vie del centro storico e proposta automatizzata di percorsi alternativi o navette dedicate a causa di eventi paesani.
\end{itemize}
Così facendo, l’implementazione di questa piattaforma punta a rendere più efficiente il sistema di trasporto per i residenti e per i turisti, rappresentando una risorsa fondamentale per l’evoluzione tecnologica e il potenziamento dei servizi di trasporto urbano in un piccolo comune, come quello di Scicli.
\subsection*{Presentazione degli attori }
Dal punto di vista dell'ingegneria del software, la modellazione del sistema ha inizio con l'identificazione dei suoi attori. Nello specifico, per la piattaforma in esame sono stati individuati i seguenti attori di sistema:
\begin{itemize}
    \item \textbf{UTENTE REGISTRATO:} rappresenta l’utente che usa il software dopo essersi registrato ed ha accesso a più funzionalità e servizi, può sia accedere alle funzionalità dell’utente non registrato sia ad altre funzionalità aggiuntive.
    \item \textbf{UTENTE NON REGISTRATO:} rappresenta l’utente che usufruisce dei servizi di base. Può consultare orari, percorsi e disponibilità dei mezzi in tempo reale, senza tuttavia poter effettuare pagamenti, prenotare un trasporto su richiesta o salvare preferenze.
    \item \textbf{TURISTA:} rappresenta un’evoluzione dell’utente registrato che ha accesso a pacchetti speciali dedicati a servizi turistici integrati, non disponibili per l’utente Registrato standard.
    \item \textbf{ADMIN DELL'ENTE:} Gestore responsabile dell’azienda di bus, o proprietari di navette turistiche. Ha la responsabilità di configurare le corse, visualizzare i guadagni e monitorare i mezzi in tempo reale.
    \item \textbf{OPERATORE:} rappresenta il profilo tecnico - operativo di colui che usa la piattaforma per gestione logistica, ricevendo in tempo reale le richieste delle chiamate del DRT, o dei taxi per vedere il percorso assegnato.
    \item \textbf{FUNZIONARIO COMUNALE:} attore istituzionale responsabile dell’aggiornamento in tempo reale di ordinanze di chiusura strade per eventi, per l’aggiornamento del traffico urbano.
    \item \textbf{PAYMENT:} attore non umano, esterno al sistema che si occupa della gestione di transazioni finanziarie.
\end{itemize}
Nelle pagine successive verranno approfonditi i diversi casi d’uso e la relativa documentazione degli attori precedentemente citati.